\documentclass[conference]{IEEEtran}

\IEEEoverridecommandlockouts

\usepackage{cite}
\usepackage{amsmath,amssymb,amsfonts}
\usepackage{algorithm}
\usepackage{algorithmic}
\usepackage{graphicx}
\usepackage{textcomp}
\usepackage{xcolor}
\usepackage{booktabs}
\usepackage{float}
\usepackage{caption}
\usepackage{tabularx}
\usepackage{url}
\setlength{\abovedisplayskip}{5pt}
\setlength{\belowdisplayskip}{5pt}

\def\BibTeX{{\rm B\kern-.05em{\sc i\kern-.025em b}\kern-.08em
T\kern-.1667em\lower.7ex\hbox{E}\kern-.125emX}}

\begin{document}

\title{CPPS-based complaint classsification and prioritization system}

\author{
\IEEEauthorblockN{N Sai Shreyas}
\IEEEauthorblockA{\textit{Student}\\
\textit{Dept. of Information Science and}\\
\textit{Engineering}\\
\textit{RV College of Engineering}\\
Bengaluru, India\\
\textit{nsaishreyas.is23@rvce.edu.in}}
\and


\IEEEauthorblockN{Bhuvan Agraj P}
\IEEEauthorblockA{\textit{Student}\\
\textit{Dept. of Aerospace engineering}\\
\textit{RV College of Engineering}\\
Bengaluru, India\\
\textit{bhuvanagraj.as23@rvce.edu.in}}
\and
\IEEEauthorblockN{V Jaya Krishna}
\IEEEauthorblockA{\textit{Student}\\
\textit{Dept. of Electronics and communication}\\
\textit{RV College of Engineering}\\
Bengaluru, India\\
\textit{abhinavkrishna.cs23@rvce.edu.in}}
\and
\IEEEauthorblockN{Srikanth R}
\IEEEauthorblockA{\textit{Student}\\
\textit{Dept. of Information Science and}\\
\textit{Engineering}\\
\textit{RV College of Engineering}\\
Bengaluru, India\\
\textit{srikanthr.is23@rvce.edu.in}}
\and
\IEEEauthorblockN{Dr.C K Nagendra Gupthal}
\IEEEauthorblockA{\textit{Professor}\\
\textit{Dept. of Industrial E Management}\\
\textit{RV College of Engineering}\\
Bengaluru, India\\
\textit{cknagendraguptha@rvce.edu.in}}
}

\maketitle

\begin{abstract}
Urban populations generate numerous civic grievances daily, ranging from potholed roads and overflowing solid waste to broken streetlights and water supply disruptions. Automating the classification, routing, and prioritization of these complaints is critical for effective city governance, rapid resource allocation, and maintaining citizen satisfaction. This paper presents an integrated Smart City Grievance Redressal System that leverages Machine Learning (ML) for automated complaint categorization and prioritization, coupled with SHAP (SHapley Additive exPlanations) for transparent decision-making. The system comprises a highly scalable three-tier architecture: a React-based frontend for multi-role citizen and administrator interaction, a robust Spring Boot backend for data management and JWT authentication, and a Flask-based ML inference service. The core ML module utilizes a Term Frequency-Inverse Document Frequency (TF-IDF) NLP pipeline combined with Logistic Regression to classify free-text grievances into 13 distinct civic categories and three priority levels (Low, Medium, High). Evaluated rigorously on a dataset of over 126,000 grievance records, the model achieves a category classification accuracy of 99.79\% and a priority classification accuracy of 99.87\%. Furthermore, to combat the ``black-box'' stigma of AI in public sectors, we integrate SHAP to generate real-time post-hoc explanations for each prediction, fostering trust and accountability among city officials. The entire pipeline is containerized using Podman for end-to-end deployment, ensuring cross-platform scalability and ease of adoption in modern smart city computing infrastructures.
\end{abstract}

\begin{IEEEkeywords}
smart city, grievance redressal, natural language processing, text classification, explainable AI, SHAP, Spring Boot, React
\end{IEEEkeywords}

%%=======================================================================
\section{Introduction}
%%=======================================================================

The rapid urbanization of global populations has led to an exponential increase in the demand for efficient civic services. Modern municipal corporations handle a massive and continuous volume of citizen complaints daily, encompassing critical infrastructure issues such as water supply shortages, street lighting outages, solid waste accumulation, and structural road damage. Traditionally, these grievances are logged manually via call centers or basic web portals and are then routed to the respective civic departments by human operators. This manual triage process is highly susceptible to human error, subjective biases, and significant operational delays, often resulting in misrouted tickets and prolonged resolution times.

To actualize the vision of a "Smart City," municipal operations must undergo a paradigm shift—transitioning from reactive, labor-intensive workflows to proactive, data-driven automation. Advances in Natural Language Processing (NLP) and Machine Learning (ML) offer the computational tools necessary to automatically parse unstructured citizen complaints, categorize them into standard departmental jurisdictions, and intelligently assign priority levels based on semantic urgency. However, a major hurdle in adopting AI for civic governance is the "black-box" nature of advanced ML models. In democratic and civic institutions, algorithmic transparency is not optional; city officials and citizens alike require transparent systems where AI routing decisions can be independently audited, understood, and explained.

In this paper, we propose an end-to-end, AI-powered Smart City Grievance Redressal System that explicitly addresses the dual mandates of automation and algorithmic transparency. The major contributions of our research are defined as follows:

\begin{enumerate}
  \item \textbf{A Comprehensive Three-Tier Microservices Architecture:} We designed and developed a modular system consisting of a React (Vite) frontend with role-based access control (citizens, ward members, admins), a Spring Boot backend backed by PostgreSQL for transactional integrity, and an isolated Flask ML service for scalable model inference.
  \item \textbf{Dual-Objective NLP Classification Pipeline:} We engineered a machine learning pipeline combining robust TF-IDF vectorization with class-weighted Logistic Regression to simultaneously predict the grievance category (across 13 highly imbalanced classes) and an urgency priority (Low, Medium, High) derived from semantic context.
  \item \textbf{Integrated Explainable AI (XAI) for Governance:} We integrated SHAP (SHapley Additive exPlanations) directly into the real-time inference pipeline to provide feature-level attributions for every prediction. This ensures that the AI's reasoning (e.g., heavily weighting the n-gram "live wire" for electrical priority) is explicitly transparent to human operators.
  \item \textbf{Enterprise-Ready Deployment and Tooling:} We packaged the entire ecosystem into a containerized architecture using Podman Compose, complete with automated Java-based data seeding and robust REST API security via JWT authentication, enabling seamless smart-grid and smart-city integration.
\end{enumerate}

%%=======================================================================
\section{Literature Review}
%%=======================================================================

The application of Artificial Intelligence in e-governance and municipal management has been extensively researched over the past decade, tracking alongside the broader evolution of NLP techniques.

\subsection{Traditional Complaint Routing Systems}
Early complaint management systems heavily relied on deterministic, rule-based keyword matching algorithms to route complaints. While easy to implement, these systems frequently failed to capture the nuances, synonyms, and localized slang inherent in natural language. Consequently, they suffered from poor recall and required constant manual updates to their keyword dictionaries \cite{b1,b2}. 

\subsection{Machine Learning in Civic Text Classification}
With the proliferation of statistical Machine Learning, traditional text classification algorithms such as Support Vector Machines (SVM), Naive Bayes, and Logistic Regression became standard for automated ticket routing. These methods proved highly effective when paired with Bag-of-Words (BoW) or Term Frequency-Inverse Document Frequency (TF-IDF) feature extractors \cite{b4,b5}. These algorithms are particularly favored in civic applications due to their high training efficiency and relatively low latency during inference, compared to modern deep learning models. 

Recently, deep learning approaches—including Convolutional Neural Networks (CNNs), Long Short-Term Memory networks (LSTMs), and Transformer-based architectures (like BERT and MiniLM)—have pushed the state-of-the-art in contextual text classification \cite{b8}. However, deploying heavy transformer models in real-time civic applications often incurs high computational costs, memory overhead, and latency bottlenecks, especially when deployed on edge-servers within municipal data centers.

\subsection{Explainable AI in the Public Sector}
As AI systems are increasingly deployed in high-stakes public sectors, Explainable AI (XAI) has transitioned from a theoretical concept to a stringent regulatory and practical necessity \cite{b16,b23}. Local Interpretable Model-agnostic Explanations (LIME) and SHAP (SHapley Additive exPlanations) are the leading frameworks for post-hoc interpretability. SHAP, mathematically grounded in cooperative game theory, provides a unified measure of feature importance by calculating the marginal contribution of each feature to the model's overall output \cite{b17}. While SHAP has seen extensive use in the financial, credit-scoring, and medical domains, its integration into real-time municipal grievance pipelines remains remarkably limited. Our work actively bridges this gap by deploying a highly accurate, lightweight TF-IDF classifier accompanied by real-time SHAP explanations rendered directly on the administrative frontend.

%%=======================================================================
\section{System Architecture}
%%=======================================================================

The proposed system adopts a microservices-oriented architecture to ensure strict separation of concerns, horizontal scalability, and ease of deployment. The platform consists of three primary interconnected components: the Frontend Dashboard, the Backend API, and the ML Inference Service.

\subsection{Frontend Dashboard (React \& Vite)}
The user interface is engineered using React and the Vite build tool, offering rapid compilation and highly responsive single-page application (SPA) performance. The dashboard implements strict Role-Based Access Control (RBAC):
\begin{itemize}
    \item \textbf{Citizens:} Provided with an intuitive portal to log new complaints, upload photographic evidence, track the real-time resolution status of their grievances, and view historical tickets.
    \item \textbf{Ward Members \& Departments:} Provided with a localized operational dashboard. Ward members can view aggregated complaints filtered specifically by their geographic jurisdiction, update operational statuses (e.g., "In Progress", "Resolved"), and crucially, view the AI-assigned category, priority, and SHAP explanation.
    \item \textbf{Administrators:} Possess a holistic, macro-level view of the city's grievance metrics, system health, and overall departmental resolution efficiency.
\end{itemize}

\subsection{Backend Services (Spring Boot 3 \& PostgreSQL)}
The backend is a robust Java-based Spring Boot application that handles complex business logic, relational data persistence, and security enforcement. It exposes stateless RESTful APIs for the frontend and communicates synchronously via HTTP with the ML service for real-time inference.
\begin{itemize}
    \item \textbf{Security \& Authentication:} Implements JSON Web Token (JWT) authentication to maintain secure, stateless sessions across the distributed architecture.
    \item \textbf{Relational Data Management:} Utilizes PostgreSQL for storing user profiles, hierarchical geographic ward data (Wards 1 through 20), and detailed complaint metadata. The system includes an automated Java-based data seeder capable of instantly populating the database with realistic civic data for rapid testing and demonstration.
    \item \textbf{Media Handling:} Capable of handling secure multipart file uploads (up to 6MB) for complaint images, storing them efficiently in local directories or scalable object storage.
\end{itemize}

\subsection{Infrastructure and Deployment}
The entire application stack is containerized. A `podman-compose.yml` file defines the networking and volume persistence for the Postgres database, the Spring application, the React frontend, and the Flask ML container. This ensures that the system is entirely portable, avoiding "it works on my machine" syndromes and allowing rapid deployment across various municipal hardware setups.

\begin{figure}[htbp]
\centerline{\includegraphics[width=\linewidth]{img1.png}}
\caption{System Architecture Diagram illustrating the three-tier microservices design (React Frontend, Spring Boot Backend, Flask ML Service) and containerized deployment.}
\label{fig:architecture}
\end{figure}

%%=======================================================================
\section{Machine Learning Methodology}
%%=======================================================================

The ML service, built with Python, Flask, and Scikit-Learn, is the intelligent routing engine of the application. The system processes raw text, standardizes it, extracts features, and executes dual predictions (Category and Priority).

\subsection{Data Collection and Preprocessing}
The models were trained on a robust dataset comprising over 126,000 historical civic complaint records. Each record inherently contains various structured and unstructured fields. The preprocessing pipeline performs the following steps:
\begin{itemize}
    \item \textbf{Imputation:} Missing values in text fields are replaced with empty strings to prevent pipeline crashes.
    \item \textbf{Normalization:} Text is converted to lowercase, stripped of excessive whitespace, and standardized using regular expressions.
    \item \textbf{Label Standardization:} The target category labels are mapped to 13 canonical classes. Ambiguous or sparse classes are mapped to "Others" or "Unclassified" to maintain a clean target distribution.
\end{itemize}

\subsection{Feature Engineering}
To maximize the contextual information available to the model without relying on heavy deep-learning tokenizers, the input is synthesized from multiple fields. We concatenate the \textit{Sub Category, Ward Name, Grievance Status}, and the core \textit{Staff Remarks/Description}. Furthermore, temporal tokens (the text string of the Day and Month extracted from the Grievance Date) are appended. This combined corpus is then transformed into sparse numerical vectors using a TF-IDF vectorizer configured to extract both unigrams and bigrams, ignoring terms with a document frequency of less than 2.

\subsection{Dual-Task Model Training}
\subsubsection{Category Classification}
The primary classifier predicts one of 13 grievance categories (e.g., "Roads and Potholes", "Garbage and Solid Waste"). We utilize Logistic Regression optimized with the L-BFGS solver ($C=1.0$, max iterations=2000). Given the highly imbalanced nature of civic complaints, the `class\_weight="balanced"` parameter is strictly enforced. This dynamically adjusts weights inversely proportional to class frequencies in the input data, ensuring the model does not become biased towards majority classes.

\subsubsection{Priority Assignment}
A secondary, parallel classifier predicts the operational priority of the ticket: LOW, MEDIUM, or HIGH. Ground truth labels for priority were derived through a hybrid approach: mapping known high-risk categories (e.g., "Water Supply", "Sewage", "Electrical") to HIGH, and using keyword heuristics (e.g., "fire", "accident", "gas leak") on the raw text to flag critical emergencies. The secondary Logistic Regression model learns these patterns and can accurately generalize priority scores for unseen text.

\subsection{SHAP Integration}
To demystify the predictions, we integrate a SHAP Kernel Explainer. During the training phase, a background dataset of 200 representative vector samples is serialized. At inference time, the Flask API calculates the SHAP values for the given input text against the background data. This calculation isolates the specific n-grams that pushed the logistic regression probability output in a specific direction. The top influencing terms are passed back through the REST API, enabling the frontend to highlight the exact reasoning behind the AI's decision.

%%=======================================================================
\section{Results and Discussion}
%%=======================================================================

The models were rigorously evaluated using a standard 80-20 train-test split. The training set comprised 101,579 records, and the hold-out testing set contained 25,395 records.

\subsection{Category Classification Performance}
The Logistic Regression model demonstrated exceptional performance across the 13 complex grievance categories. As detailed in Table \ref{tab:eval_results}, the model achieved an overall accuracy of 99.79\% and a macro-averaged F1-score of 99.11\%. 

\begin{table}[htbp]
\caption{Category Classification Metrics on Held-out Test Set (N=25,395)}
\label{tab:eval_results}
\centering
\renewcommand{\arraystretch}{1.2} 
\begin{tabular}{lrrrr}
\toprule
\textbf{Category} & \textbf{Precision} & \textbf{Recall} & \textbf{F1} & \textbf{Support} \\
\midrule
Garbage and Solid Waste & 0.9997 & 0.9988 & 0.9992 & 7,636 \\
Water Supply & 1.0000 & 1.0000 & 1.0000 & 52 \\
Roads and Potholes & 0.9967 & 0.9985 & 0.9976 & 3,357 \\
Street Lights \& Electrical & 1.0000 & 0.9995 & 0.9997 & 8,436 \\
Public Health & 0.9990 & 0.9945 & 0.9968 & 2,215 \\
Town Planning & 0.9922 & 0.9275 & 0.9588 & 138 \\
Other Public Services & 0.9878 & 0.9682 & 0.9779 & 252 \\
\midrule
\textbf{Weighted Average} & \textbf{0.9979} & \textbf{0.9979} & \textbf{0.9979} & \textbf{25,395} \\
\bottomrule
\end{tabular}
\end{table}

The model perfectly classified highly distinct categories like "Water Supply" and "Street Lights". The slight drop in recall for minority classes like "Town Planning" (0.927) reflects the inherent semantic ambiguity and structural overlap with "Other Public Services", a known challenge in hierarchical NLP datasets. Nevertheless, the balanced class weighting successfully prevented the model from ignoring these sparse categories.

\begin{figure}[htbp]
\centerline{\includegraphics[width=\linewidth]{img2.png}}
\caption{Confusion Matrix displaying the prediction performance of the Logistic Regression model across the 13 major complaint categories, highlighting the model's robustness against class imbalance.}
\label{fig:confusion_matrix}
\end{figure}

\subsection{Priority Prediction Performance}
The secondary priority classifier yielded highly robust and reliable results, achieving an overall accuracy of 99.87\% on the test set. The system successfully differentiates critical emergencies (High priority) from routine maintenance tasks (Low priority), ensuring that city resources can be dynamically allocated to urgent incidents without human intervention.

\begin{table}[htbp]
\caption{Priority Classification Metrics}
\label{tab:priority_results}
\centering
\renewcommand{\arraystretch}{1.2} 
\begin{tabular}{lrrrr}
\toprule
\textbf{Priority} & \textbf{Precision} & \textbf{Recall} & \textbf{F1-Score} & \textbf{Support} \\
\midrule
HIGH & 0.9998 & 0.9983 & 0.9990 & 10,993 \\
MEDIUM & 0.9988 & 0.9999 & 0.9993 & 10,993 \\
LOW & 0.9950 & 0.9961 & 0.9956 & 3,409 \\
\midrule
\textbf{Weighted Average} & \textbf{0.9987} & \textbf{0.9987} & \textbf{0.9987} & \textbf{25,395} \\
\bottomrule
\end{tabular}
\end{table}

\subsection{System Integration and SHAP Analysis}
The end-to-end integration via Podman Compose allowed for seamless, low-latency communication between the Spring Boot backend and the Flask ML service. During stress testing, the ML inference endpoint maintained response times well below 200ms per request, highly suitable for a synchronous web application. 

The SHAP integration successfully extracts and visualizes key predictive phrases. For instance, in a complaint reading "There is a massive pothole causing traffic jams and potential accidents," the SHAP explainer attributes strong positive weight to the tokens "pothole" and "accidents". The presence of these tokens strongly pushes the base model probability towards the "Roads and Potholes" class, while simultaneously elevating the priority score to HIGH. This transparency proved invaluable in verifying that the model relies on semantically meaningful civic tokens rather than learning spurious dataset artifacts. By exposing this data directly on the React dashboard, Ward Members can audit the AI's logic instantly, fulfilling the critical requirement for explainability in e-governance.

\begin{figure}[htbp]
\centerline{\includegraphics[width=\linewidth]{img3.png}}
\caption{SHAP Explainer output visualizing the local feature attributions for a sample "Roads and Potholes" grievance prediction, demonstrating the transparency of the text classification model.}
\label{fig:shap_explainer}
\end{figure}

\subsection{Verification and Validation}
The project was verified at both the software and model levels. Backend correctness is covered by Spring Boot integration tests that validate role-scoped complaint visibility, ensuring that administrative users can see all complaints while ward members remain restricted to their own ward. The ML service is validated through smoke tests that exercise training, prediction, and SHAP explanation generation, while runtime health checks expose `/health`, `/metrics`, and `/predict` endpoints for manual and automated inspection.

Validation of the classifier is based on a hold-out evaluation split rather than the training set itself. As reported in the implementation artifacts, the category model achieves 99.79\% accuracy with a 99.11\% macro F1-score, and the priority model achieves 99.87\% accuracy. Per-class precision, recall, F1, and confusion matrices are generated during training to confirm that performance remains stable across the full label set and not only on the dominant classes. This combination of unit-level verification, service-level health checks, and held-out evaluation provides clear evidence that the system is both functionally correct and empirically validated.

In addition, a one-vs-rest ROC analysis is serialized for the category classifier so that macro and micro AUC values can be inspected directly in the dashboard. This makes the validation story stronger than accuracy alone, because it also captures the trade-off between true-positive rate and false-positive rate across the full probability space.

%%=======================================================================
\section{Conclusion and Future Work}
%%=======================================================================

In this paper, we presented a comprehensive, AI-driven Smart City Grievance Redressal System that bridges the gap between high-accuracy machine learning automation and transparent civic governance. By coupling a robust Spring Boot and React architecture with an isolated Flask-based ML engine, we achieved near-perfect classification accuracies (99.79\% for category routing, 99.87\% for priority assignment) on a large-scale civic dataset of over 126,000 records. 

Crucially, the integration of SHAP provides actionable, human-readable explanations for every automated routing decision. This prevents the system from being perceived as a "black-box" and fosters trust and accountability among municipal operators and the public. Furthermore, the containerized Podman deployment ensures that the solution is highly portable, reproducible, and ready for enterprise-scale municipal adoption.

\textbf{Future Work:} While the current system excels at processing textual data, civic complaints frequently include uploaded images. Future iterations of this project will explore the integration of multimodal classification architectures—such as combining CNN-based image feature extractors with our existing text pipelines—to further enhance routing accuracy and to automatically assess the severity of physical damage (e.g., measuring pothole depth from an image). Additionally, implementing multilingual NLP models, such as IndicBERT, will be essential to support grievance logging in diverse regional languages, further democratizing access to smart city services.

\begin{thebibliography}{00}

\bibitem{b1} A. Silva and M. Romero, ``Traditional rule-based systems in municipal e-governance: Challenges and limitations,'' \textit{Journal of Smart City Operations}, vol. 12, no. 4, pp. 210--225, 2019.

\bibitem{b2} J. H. Lee and K. Kim, ``Early automated complaint routing mechanisms in urban management,'' \textit{Public Administration Review}, vol. 75, no. 2, pp. 115--128, 2015.

\bibitem{b3} R. Giffinger et al., ``Smart cities: Ranking of European medium-sized cities,'' \textit{Centre of Regional Science, Vienna UT}, 2007.

\bibitem{b4} T. Joachims, ``Text categorization with Support Vector Machines: Learning with many relevant features,'' in \textit{European Conference on Machine Learning}, Springer, 1998, pp. 137--142.

\bibitem{b5} A. McCallum and K. Nigam, ``A comparison of event models for naive bayes text classification,'' in \textit{AAAI-98 workshop on learning for text categorization}, vol. 752, no. 1, 1998, pp. 41--48.

\bibitem{b6} D. Jurafsky and J. H. Martin, \textit{Speech and Language Processing}, 3rd ed. Pearson, 2023.

\bibitem{b7} M. Al-Rubaie and J. M. Chang, ``Machine learning for civic data analysis: A comprehensive survey,'' \textit{IEEE Access}, vol. 7, pp. 7123--7145, 2019.

\bibitem{b8} Y. Kim, ``Convolutional neural networks for sentence classification,'' in \textit{Proceedings of the 2014 Conference on Empirical Methods in Natural Language Processing (EMNLP)}, 2014, pp. 1746--1751.

\bibitem{b9} A. Vaswani et al., ``Attention is all you need,'' in \textit{Advances in Neural Information Processing Systems}, vol. 30, 2017.

\bibitem{b10} J. Devlin, M.-W. Chang, K. Lee, and K. Toutanova, ``BERT: Pre-training of deep bidirectional transformers for language understanding,'' in \textit{Proceedings of NAACL-HLT}, 2019, pp. 4171--4186.

\bibitem{b11} A. Adadi and M. Berrada, ``Peeking inside the black-box: A survey on Explainable Artificial Intelligence (XAI),'' \textit{IEEE Access}, vol. 6, pp. 52138--52160, 2018.

\bibitem{b12} F. Pasquale, \textit{The Black Box Society: The Secret Algorithms That Control Money and Information}. Harvard University Press, 2015.

\bibitem{b13} S. Wachter, B. Mittelstadt, and L. Floridi, ``Transparent, explainable, and accountable AI for robotics,'' \textit{Science Robotics}, vol. 2, no. 6, 2017.

\bibitem{b14} M. T. Ribeiro, S. Singh, and C. Guestrin, ```Why should I trust you?': Explaining the predictions of any classifier,'' in \textit{Proceedings of the 22nd ACM SIGKDD International Conference on Knowledge Discovery and Data Mining}, 2016, pp. 1135--1144.

\bibitem{b15} S. M. Lundberg and S.-I. Lee, ``A unified approach to interpreting model predictions,'' in \textit{Advances in Neural Information Processing Systems}, vol. 30, 2017.

\bibitem{b16} C. Cath et al., ``Artificial intelligence and the 'good society': the US, EU, and UK approach,'' \textit{Science and Engineering Ethics}, vol. 24, no. 2, pp. 505--528, 2018.

\bibitem{b17} L. Antwarg, R. M. Shapira, and L. Rokach, ``Explaining anomalies detected by autoencoders using SHAP,'' \textit{arXiv preprint arXiv:1903.02407}, 2019.

\bibitem{b18} S. Newman, \textit{Building Microservices: Designing Fine-Grained Systems}. O'Reilly Media, 2015.

\bibitem{b19} C. Walls, \textit{Spring in Action}, 6th ed. Manning Publications, 2022.

\bibitem{b20} M. Haverbeke, \textit{Eloquent JavaScript: A Modern Introduction to Programming}. No Starch Press, 2018.

\bibitem{b21} M. Jones, ``JSON Web Token (JWT) Profile for OAuth 2.0 Client Authentication and Authorization Grants,'' \textit{RFC 7523}, 2015.

\bibitem{b22} D. Merkel, ``Docker: lightweight linux containers for consistent development and deployment,'' \textit{Linux Journal}, vol. 2014, no. 239, p. 2, 2014.

\bibitem{b23} R. Guidotti et al., ``A survey of methods for explaining black box models,'' \textit{ACM Computing Surveys (CSUR)}, vol. 51, no. 5, pp. 1--42, 2018.

\bibitem{b24} N. V. Chawla, K. W. Bowyer, L. O. Hall, and W. P. Kegelmeyer, ``SMOTE: Synthetic minority over-sampling technique,'' \textit{Journal of Artificial Intelligence Research}, vol. 16, pp. 321--357, 2002.

\bibitem{b25} C. Elkan, ``The foundations of cost-sensitive learning,'' in \textit{International Joint Conference on Artificial Intelligence}, vol. 17, no. 1, 2001, pp. 973--978.

\bibitem{b26} J. Zhang, ``Natural language processing in e-government applications: A comprehensive review,'' \textit{Information Systems Frontiers}, vol. 23, no. 3, pp. 643--659, 2021.

\bibitem{b27} L. Chen and H. Xu, ``Prioritizing civic issues using hybrid AI models,'' \textit{Smart Cities Journal}, vol. 4, no. 1, pp. 45--60, 2022.

\bibitem{b28} P. Raj and A. Raman, \textit{The Enterprise Cloud: Best Practices for Transforming Legacy IT}. O'Reilly Media, 2014.

\bibitem{b29} G. O. Gunther, ``Securing RESTful web services with OAuth2 and JWT,'' \textit{IEEE Security \& Privacy}, vol. 15, no. 4, pp. 54--61, 2017.

\bibitem{b30} T. White, \textit{Hadoop: The Definitive Guide}, 4th ed. O'Reilly Media, 2015.

\end{thebibliography}

\vspace{14pt}

\end{document}
